% language: swedish
% figure: fig1 0.205 0.547 0.405 0.643
\section{Föreläsning 10/9}

\subsection{Komplexa logaritmer}
\begin{definition}
Givet område $G$ så sägs varje \emph{kontinuerlig} funktion $\operatorname{Log} z : G \to \mathbb{C}$ som uppfyller $e^{\operatorname{Log} z} = z$ vara en \emph{gren} av logaritmen $\log z$. \ $\operatorname{Log} z = \ln|z| + i\operatorname{Arg} z$ \textcolor{red!70!black}{$(-\pi < \operatorname{Arg} z \le \pi)$} kallas den \emph{principala grenen}.
\end{definition}

\begin{example}
$\operatorname{Log}(3i) = \ln 3 + i\pi/2$

$\operatorname{Log}(-7) = \ln 7 + i\pi$ \quad \textcolor{magenta!80!black}{(fick ej ta $-\pi$ så då tar vi $\pi$ istället)}
\end{example}

\begin{note}
$\operatorname{Log} z$ är diskontinuerlig längs $\mathbb{R}_{\le 0}$, så $\operatorname{Log} z$ är \emph{inte} en gren av $\log z$ på $\{\operatorname{Re} z < 0\}$
\end{note}

\begin{note}
Ibland $\operatorname{Log} z^2 \neq 2\operatorname{Log} z$

e.g.\ $\operatorname{Log}((-7)^2) = \ln 49 = 2\ln 7$

\qquad $2\operatorname{Log}(-7) = 2\ln 7 + 2\pi i$
\end{note}

\begin{example}
\notefigure{fig1}{}
Om $\operatorname{Log} z$ är en gren av $\log z$ på $G$, och $\operatorname{Log} 1 = 0$, vad blir $\operatorname{Log} 2$?

\textcolor{red!70!black}{\textbf{Svar:}} $\ln 2 - i2\pi$
\end{example}

\noindent\rule{\linewidth}{0.4pt}

\textcolor{blue!60!black}{
$\operatorname{Log} z \leftarrow$ principella \\
$\mathcal{L}og\, z$ $\leftarrow$ godtycklig \\
$\log z \leftarrow$ någon gren}
